\documentclass[11pt]{article}

\usepackage[T1]{fontenc}
\usepackage[utf8]{inputenc}
\usepackage{lmodern}
\usepackage[a4paper,margin=28mm]{geometry}
\usepackage{microtype}
\usepackage{amsmath,amssymb,amsthm,mathtools,bm}
\usepackage{booktabs,tabularx,array}
\usepackage{enumitem}
\usepackage[dvipsnames]{xcolor}
\usepackage{hyperref}
\usepackage[nameinlink,noabbrev]{cleveref}

\hypersetup{
  colorlinks=true,
  linkcolor=MidnightBlue,
  citecolor=ForestGreen,
  urlcolor=BrickRed,
  pdftitle={Gaussian Wasserstein barycenters and contraction of the raw covariance map},
  pdfauthor={Technical study prepared with Codex}
}

\setlength{\parindent}{0pt}
\setlength{\parskip}{0.55em}
\setlist[itemize]{topsep=0.25em,itemsep=0.2em,leftmargin=1.6em}
\setlist[enumerate]{topsep=0.25em,itemsep=0.2em,leftmargin=1.8em}

\newtheorem{theorem}{Theorem}[section]
\newtheorem{proposition}[theorem]{Proposition}
\newtheorem{lemma}[theorem]{Lemma}
\newtheorem{corollary}[theorem]{Corollary}
\newtheorem{conjecture}[theorem]{Conjecture}
\theoremstyle{definition}
\newtheorem{definition}[theorem]{Definition}
\newtheorem{remark}[theorem]{Remark}
\newtheorem{example}[theorem]{Example}

\newcommand{\R}{\mathbb R}
\newcommand{\Splus}{\mathbb S_+^d}
\newcommand{\Spp}{\mathbb S_{++}^d}
\newcommand{\Id}{I_d}
\newcommand{\tr}{\operatorname{tr}}
\newcommand{\diag}{\operatorname{diag}}
\newcommand{\Lip}{\operatorname{Lip}}
\newcommand{\spec}{\operatorname{spec}}
\newcommand{\rhoop}{\rho}
\newcommand{\KL}{\operatorname{KL}}
\newcommand{\cN}{\mathcal N}
\newcommand{\cP}{\mathcal P}
\newcommand{\cG}{\mathcal G}
\newcommand{\cK}{\mathcal K}
\newcommand{\cL}{\mathcal L}
\newcommand{\cT}{\mathcal T}
\newcommand{\cH}{\mathcal H}
\newcommand{\eps}{\varepsilon}
\newcommand{\dd}{\,\mathrm d}
\newcommand{\ip}[2]{\left\langle #1,#2\right\rangle}
\newcommand{\norm}[1]{\left\lVert #1\right\rVert}
\newcommand{\abs}[1]{\left\lvert #1\right\rvert}
\newcommand{\preceqL}{\preccurlyeq}

\title{\bfseries Gaussian Wasserstein Barycenters and the Raw Covariance Map\\
\large Fixed points, algorithms, contraction on compact spectral sets,\\
and entropic and unbalanced extensions}
\author{Technical study}
\date{Literature status and analysis as of 13 July 2026}

\begin{document}
\maketitle

\begin{abstract}
For positive weights $\lambda_s$ summing to one and positive definite covariance
matrices $\Sigma_s$, the centered Gaussian $2$-Wasserstein barycenter has the
unique covariance $\Sigma_\star$ satisfying
\[
  \Sigma_\star=\Psi_\lambda(\Sigma_\star),\qquad
  \Psi_\lambda(X):=\sum_s\lambda_s
  \bigl(X^{1/2}\Sigma_sX^{1/2}\bigr)^{1/2}.
\]
The fixed-point identity is classical, but it does not by itself justify the raw
Picard iteration $X_{n+1}=\Psi_\lambda(X_n)$.  This note separates that raw map
from the globally convergent transport/gradient update
\[
  K(X)=X^{-1/2}\Psi_\lambda(X)^2X^{-1/2}.
\]
It proves exact scalar and commuting results, derives the full Fr\'echet derivative
of $\Psi_\lambda$ through two Sylvester equations, gives an explicit sufficient
contraction certificate on compact spectral boxes, and formulates the remaining
noncommutative issue as a precise derivative-product problem.  In particular,
on each fixed invariant scalar interval bounded away from zero, a sufficiently
high iterate $\Psi_\lambda^k$ is a contraction; nevertheless, no fixed $k$ works
uniformly over all such intervals or all condition numbers.  Thus a common
one-dimensional objection has the wrong quantifier, while a general matrix
contraction theorem remains unsupported by the literature located for this
study.

The note also proves Gaussian preservation of ordinary Wasserstein barycenters
from Gelbrich's nonexpansive Gaussianization map, distinguishes three inequivalent
notions of entropic/Sinkhorn barycenter, and explains why the same Gaussianization
argument fails for ambient Hellinger--Kantorovich/Wasserstein--Fisher--Rao
geometry.  At the Hellinger endpoint the barycenter of Gaussian densities is
generically a Gaussian mixture, not a Gaussian.  A Gaussian-restricted BWFR
geometry is presented as a valid but explicitly constrained alternative.
\end{abstract}

\tableofcontents

\section{Executive conclusions}

The conclusions are easiest to state before the details.

\begin{enumerate}
  \item A compact subset of strictly positive definite matrices must be bounded
  away from the boundary.  The natural model domain is
  \[
    \cK_{m,M}:=\{X\in\Spp:m\Id\preceq X\preceq M\Id\},
    \qquad 0<m\leq M<\infty.
  \]
  Saying that a map is ``contractive'' is incomplete until a metric, a domain,
  and whether the map preserves that domain have been specified.

  \item The Gaussian covariance equation has a unique positive definite
  solution.  This follows from the Wasserstein barycenter theory of
  Agueh--Carlier and from matrix-analytic convexity arguments
  \cite{AguehCarlier2011,BhatiaJainLim2019}.  Brouwer's theorem proves existence
  of a fixed point on a compact spectral box; it does not prove that the raw
  Picard iterates converge.

  \item The algorithm with a classical global convergence theorem is
  \[
    X_{n+1}=K(X_n):=X_n^{-1/2}\Psi_\lambda(X_n)^2X_n^{-1/2},
  \]
  not $X_{n+1}=\Psi_\lambda(X_n)$.  The two maps have the same positive
  definite fixed point but different dynamics.  In one dimension $K$ reaches
  the answer in one step, whereas the raw map only halves the exponent at each
  iteration.

  \item In dimension one, with
  $c:=\sum_s\lambda_s\sqrt{\Sigma_s}$ and $x_\star=c^2$,
  \[
    \Psi(x)=c\sqrt{x},\qquad
    \Psi^k(x)=x_\star^{1-2^{-k}}x^{2^{-k}}.
  \]
  On an invariant interval $[m,M]$ the exact Euclidean contraction constant is
  \[
    q_k=2^{-k}\left(\frac{x_\star}{m}\right)^{1-2^{-k}}.
  \]
  Hence $q_k\to0$ for every fixed $m>0$, but for every prescribed $k$ one can
  choose $m$ so small that $q_k>1$.  The correct statement is therefore:
  eventual contraction holds on each fixed compact interval, but there is no
  universal iterate independent of the lower spectral bound.

  \item For general noncommuting matrices, this note gives (i) the exact
  derivative, (ii) a computable local spectral-radius criterion, and (iii) a
  conservative analytic compact-box certificate.  These yield rigorous
  contraction results in several regimes, including commuting data and
  sufficiently well-conditioned boxes.  A theorem covering every collection
  of positive definite inputs on every invariant compact spectral box was not
  found in the primary literature.  The older numerical reports concerning
  the raw map are not rigorous counterexamples.

  \item Gaussian Wasserstein barycenters are Gaussian because moment-matching
  Gaussianization is $1$-Lipschitz for $W_2$ by Gelbrich's inequality.  This is
  a nonexpansive projection argument, not a strict contraction argument.

  \item ``Entropic barycenter'' can mean a product-reference entropic OT
  barycenter, a Schrodinger barycenter with Lebesgue/Brownian reference, or a
  debiased Sinkhorn-divergence barycenter.  They have different covariance
  equations and opposite variance biases.  General contraction of their
  covariance Picard maps is not known; scalar derivatives already rule out a
  generic one-step Euclidean contraction near the positive-semidefinite
  boundary.

  \item Ordinary ambient HK/WFR barycenters do not admit a Gelbrich-style
  Gaussianization proof.  Moment Gaussianization is not even continuous for
  HK on a compact covariance range.  At the pure Hellinger endpoint, the exact
  barycenter of Gaussian-shaped densities is generically a finite Gaussian
  mixture.  Gaussian-restricted BWFR barycenters are meaningful finite-
  dimensional surrogates, but Gaussianity then holds by restriction, not by an
  ambient preservation theorem.
\end{enumerate}

\section{Problem and notation}

Let $\Sigma_1,\ldots,\Sigma_N\in\Spp$ and let
$\lambda_s>0$ with $\sum_{s=1}^N\lambda_s=1$.  On $\Splus$ define the
Bures--Wasserstein distance
\begin{equation}\label{eq:bures-distance}
 d_{\mathrm{BW}}^2(A,B)
 :=\tr A+\tr B-2\tr\bigl(A^{1/2}BA^{1/2}\bigr)^{1/2}.
\end{equation}
For Gaussian laws,
\begin{equation}\label{eq:gaussian-w2}
 W_2^2\bigl(\cN(m,A),\cN(n,B)\bigr)
 =\norm{m-n}_2^2+d_{\mathrm{BW}}^2(A,B).
\end{equation}
The covariance part of the barycenter problem is
\begin{equation}\label{eq:barycenter-objective}
 \min_{X\in\Splus} F(X),\qquad
 F(X):=\frac12\sum_{s=1}^N\lambda_s
 d_{\mathrm{BW}}^2(X,\Sigma_s).
\end{equation}

The map under study is
\begin{equation}\label{eq:Psi-def}
 \boxed{\quad
 \Psi(X)=\Psi_\lambda(X):=\sum_{s=1}^N\lambda_s
 R_s(X),\qquad
 R_s(X):=\bigl(X^{1/2}\Sigma_sX^{1/2}\bigr)^{1/2}.
 \quad}
\end{equation}
It is continuous and smooth on $\Spp$, maps $\Spp$ into $\Spp$, and is
positively homogeneous of degree $1/2$:
\begin{equation}\label{eq:homogeneity}
 \Psi(tX)=t^{1/2}\Psi(X),\qquad t>0.
\end{equation}
It is orthogonally equivariant, but it is not generally equivariant under an
arbitrary congruence.  That loss of full congruence equivariance is one source
of difficulty in the noncommutative contraction problem.

We use $\norm{\cdot}_F$ for the Frobenius norm and $\norm{\cdot}_{\mathrm{op}}$
for the spectral norm.  Unless explicitly stated otherwise, ``contraction''
below means a strict Lipschitz contraction for $\norm{\cdot}_F$.

\begin{remark}[Positive definite versus positive semidefinite]
The phrase ``a compact domain of positive PSD matrices'' should be made
precise.  A compact subset of $\Spp$ has a uniform lower eigenvalue bound
$m>0$.  If singular matrices are admitted and the domain meets the boundary,
the scalar square root is not Lipschitz there.  Thus no Euclidean Lipschitz
theorem uniform down to $0$ can hold even in dimension one.
\end{remark}

\section{Why Wasserstein barycenters of Gaussians are Gaussian}

\subsection{Gelbrich's inequality}

For a probability measure $\mu\in\cP_2(\R^d)$, write $m_\mu$ and
$C_\mu$ for its mean and covariance.  Its moment-matching Gaussianization is
\begin{equation}\label{eq:gaussianization}
 \cG(\mu):=\cN(m_\mu,C_\mu).
\end{equation}

\begin{theorem}[Gelbrich bound]\label{thm:gelbrich}
For $\mu,\nu\in\cP_2(\R^d)$,
\begin{equation}\label{eq:gelbrich}
 W_2^2(\mu,\nu)\geq
 \norm{m_\mu-m_\nu}_2^2+
 d_{\mathrm{BW}}^2(C_\mu,C_\nu).
\end{equation}
Equality holds for Gaussian measures and, more generally, for compatible
location-scatter families.  Consequently,
\begin{equation}\label{eq:G-nonexpansive}
 W_2\bigl(\cG(\mu),\cG(\nu)\bigr)\leq W_2(\mu,\nu).
\end{equation}
\end{theorem}

\begin{proof}
Take any coupling $(X,Y)$ of $\mu$ and $\nu$ and center it:
$\widetilde X=X-m_\mu$, $\widetilde Y=Y-m_\nu$.  If
$C=\mathbb E[\widetilde X\widetilde Y^\top]$, then
\begin{equation}\label{eq:coupling-cost-cross}
 \mathbb E\norm{X-Y}_2^2
 =\norm{m_\mu-m_\nu}_2^2+
 \tr C_\mu+\tr C_\nu-2\tr C.
\end{equation}
The block covariance is positive semidefinite:
\[
 \begin{pmatrix}C_\mu&C\\C^\top&C_\nu\end{pmatrix}\succeq0.
\]
When the covariances are invertible, the block factorization criterion gives
$C=C_\mu^{1/2}KC_\nu^{1/2}$ for a contraction $K$; the singular case follows
by approximation.  Von Neumann's trace inequality, or equivalently the polar
factor maximization, yields
\[
 \tr C\leq
 \tr\bigl(C_\mu^{1/2}C_\nu C_\mu^{1/2}\bigr)^{1/2}.
\]
Substitution in \eqref{eq:coupling-cost-cross} and minimization over couplings
give \eqref{eq:gelbrich}.  For two Gaussians, the maximizing cross-covariance
is realized by a jointly Gaussian coupling, which proves equality and
\eqref{eq:G-nonexpansive}.
\end{proof}

This result is due to Gelbrich \cite{Gelbrich1990}; the Gaussian distance
formula was obtained earlier in
\cite{DowsonLandau1982,OlkinPukelsheim1982,GivensShortt1984}.

\subsection{Gaussian preservation of the barycenter}

\begin{theorem}[Gaussian barycenter by Gaussianization]\label{thm:gaussian-barycenter}
Let $\nu_s=\cN(m_s,\Sigma_s)$ with $\Sigma_s\succ0$.  The unique minimizer of
\begin{equation}\label{eq:measure-barycenter}
 \mu\longmapsto \frac12\sum_s\lambda_sW_2^2(\mu,\nu_s)
\end{equation}
over $\cP_2(\R^d)$ is Gaussian.  Its mean is
$m_\star=\sum_s\lambda_sm_s$, and its covariance is the unique minimizer of
\eqref{eq:barycenter-objective}.
\end{theorem}

\begin{proof}
Since $\cG(\nu_s)=\nu_s$, \cref{thm:gelbrich} gives, for every candidate
$\mu$,
\[
 W_2\bigl(\cG(\mu),\nu_s\bigr)\leq W_2(\mu,\nu_s)
 \quad\text{for every }s.
\]
Thus Gaussianization cannot increase the barycenter objective.  General
Wasserstein barycenter theory gives uniqueness when at least one input is
absolutely continuous; here every nondegenerate Gaussian is absolutely
continuous.  Hence the unique minimizer equals its Gaussianization.  Formula
\eqref{eq:gaussian-w2} separates mean and covariance, and the mean minimizer is
the weighted Euclidean mean.
\end{proof}

\begin{remark}
The Gaussianization map is nonexpansive, not strictly contractive.  Equality
holds for every pair of Gaussians.  Gaussian preservation should therefore not
be confused with a Banach contraction mechanism.
\end{remark}

\section{The covariance equation and its unique solution}

\subsection{Optimal maps and first-order condition}

The optimal linear transport from $\cN(0,X)$ to $\cN(0,\Sigma_s)$ is
\begin{equation}\label{eq:optimal-map}
 T_s(X):=X^{-1/2}R_s(X)X^{-1/2}.
\end{equation}
It is symmetric positive definite and satisfies
$T_s(X)XT_s(X)=\Sigma_s$.  Define its weighted mean
\begin{equation}\label{eq:H-def}
 H(X):=\sum_s\lambda_sT_s(X).
\end{equation}
Since $R_s(X)=X^{1/2}T_s(X)X^{1/2}$,
\begin{equation}\label{eq:Psi-H}
 \Psi(X)=X^{1/2}H(X)X^{1/2}.
\end{equation}

\begin{proposition}[Differential of the Bures objective]
For a symmetric perturbation $Z$,
\begin{equation}\label{eq:bures-differential}
 D\,d_{\mathrm{BW}}^2(X,\Sigma_s)[Z]
 =\tr\bigl((\Id-T_s(X))Z\bigr),
\end{equation}
and therefore
\begin{equation}\label{eq:F-gradient}
 \nabla F(X)=\frac12(\Id-H(X)).
\end{equation}
\end{proposition}

\begin{proof}
The non-linear trace term may equivalently be written
\[
 \tr R_s(X)=
 \tr\bigl(\Sigma_s^{1/2}X\Sigma_s^{1/2}\bigr)^{1/2}.
\]
For a positive definite $Y$, spectral calculus and cyclicity of the trace give
\[
 D\,\tr(Y^{1/2})[W]=\frac12\tr(Y^{-1/2}W).
\]
With $Y=\Sigma_s^{1/2}X\Sigma_s^{1/2}$, this yields
\[
 D\,\tr R_s(X)[Z]
 =\frac12\tr\left(
 \Sigma_s^{1/2}(\Sigma_s^{1/2}X\Sigma_s^{1/2})^{-1/2}
 \Sigma_s^{1/2}Z\right)
 =\frac12\tr(T_s(X)Z),
\]
where the last matrix is the equivalent target-side formula for the Gaussian
optimal map.  Differentiating \eqref{eq:bures-distance} proves
\eqref{eq:bures-differential}; summing proves \eqref{eq:F-gradient}.
\end{proof}

Thus the stationarity equation is
\begin{equation}\label{eq:stationarity-H}
 H(X)=\Id.
\end{equation}
By \eqref{eq:Psi-H}, this is equivalent to
\begin{equation}\label{eq:fixed-equation}
 \boxed{\Psi(X)=X.}
\end{equation}

\begin{theorem}[Existence and uniqueness]\label{thm:exist-unique}
For $\Sigma_s\succ0$ and positive weights summing to one, there is a unique
$\Sigma_\star\succ0$ satisfying \eqref{eq:fixed-equation}.  It is the unique
minimizer of \eqref{eq:barycenter-objective}.
\end{theorem}

This is the Gaussian part of Agueh--Carlier's barycenter theorem
\cite[Theorem 6.1]{AguehCarlier2011}.  Earlier work derived the equation and
studied the associated multi-coupling problem
\cite{KnottSmith1994,RuschendorfUckelmann2002}.  Bhatia--Jain--Lim give a
matrix-analysis proof based on convexity of the fidelity functional
\cite{BhatiaJainLim2019}.

For completeness, strict uniqueness can also be seen directly.  The map
$Y\mapsto\tr(Y^{1/2})$ is strictly concave on positive definite matrices.  As
$X\mapsto\Sigma_s^{1/2}X\Sigma_s^{1/2}$ is injective when
$\Sigma_s\succ0$, each function
\[
 X\longmapsto \tr X+\tr\Sigma_s
 -2\tr(\Sigma_s^{1/2}X\Sigma_s^{1/2})^{1/2}
\]
is strictly convex.  Hence $F$ has at most one stationary point, and the
fixed point constructed on the compact box below is that unique minimizer.

\subsection{The compact spectral box used for existence}

Assume
\begin{equation}\label{eq:data-bounds}
 a_s\Id\preceq\Sigma_s\preceq b_s\Id,
 \qquad 0<a_s\leq b_s.
\end{equation}
Set
\begin{equation}\label{eq:cpm}
 c_-:=\sum_s\lambda_s\sqrt{a_s},\qquad
 c_+:=\sum_s\lambda_s\sqrt{b_s}.
\end{equation}

\begin{lemma}[Spectral enclosure]\label{lem:spectral-enclosure}
If $X\in\cK_{m,M}$, then
\begin{equation}\label{eq:image-box}
 c_-\sqrt m\,\Id\preceq\Psi(X)
 \preceq c_+\sqrt M\,\Id.
\end{equation}
Consequently, $\cK_{m,M}$ is forward invariant under $\Psi$ whenever
\begin{equation}\label{eq:invariance-condition}
 m\leq c_-^2,\qquad M\geq c_+^2.
\end{equation}
\end{lemma}

\begin{proof}
From \eqref{eq:data-bounds},
\[
 a_sX\preceq X^{1/2}\Sigma_sX^{1/2}\preceq b_sX.
\]
The matrix square root is operator monotone, so for $X\in\cK_{m,M}$,
\[
 \sqrt{a_sm}\,\Id\preceq R_s(X)
 \preceq\sqrt{b_sM}\,\Id.
\]
Sum with weights.  Conditions \eqref{eq:invariance-condition} are exactly
$c_-\sqrt m\geq m$ and $c_+\sqrt M\leq M$.
\end{proof}

Thus $\Psi$ is a continuous self-map of a convex compact set, and Brouwer
gives a fixed point.  This is an existence argument only.  A continuous
self-map of a compact convex set need not be a contraction, and its Picard
iterates need not converge.

\section{Algorithms in the literature: the essential distinction}

\subsection{Raw Picard versus the transport update}

The raw fixed-point iteration is
\begin{equation}\label{eq:raw-picard}
 X_{n+1}=\Psi(X_n).
\end{equation}
The measure pushforward construction of Alvarez-Esteban et al. instead gives
the covariance update
\begin{equation}\label{eq:K-update}
 \boxed{\quad
 K(X):=H(X)XH(X)
 =X^{-1/2}\Psi(X)^2X^{-1/2}.
 \quad}
\end{equation}
The identity follows from \eqref{eq:Psi-H}.  Both maps have the same positive
definite fixed point, but they are not the same map.

More generally, a damped Bures--Wasserstein gradient step has the form
\begin{equation}\label{eq:damped-rgd}
 K_\eta(X)=\bigl((1-\eta)\Id+\eta H(X)\bigr)
 X\bigl((1-\eta)\Id+\eta H(X)\bigr),
\end{equation}
with normalization depending on whether the objective contains a factor
$1/2$.  The unit step $\eta=1$ is \eqref{eq:K-update}.

\begin{theorem}[Global convergence of the stabilized update]
Suppose the input covariances are positive semidefinite and at least one is
positive definite.  Starting from $X_0\succ0$, the sequence
$X_{n+1}=K(X_n)$ converges to the Bures--Wasserstein barycenter covariance.
Moreover, the barycenter variance decreases and controls the step:
\begin{equation}\label{eq:variance-descent}
 F(X)\geq F(K(X))+\frac12d_{\mathrm{BW}}^2(X,K(X)),
\end{equation}
up to the normalization of $F$.
\end{theorem}

This is the covariance specialization of
\cite[Theorem 4.2]{AlvarezEsteban2016}; see also
\cite{BhatiaJainLim2019}.  Determinant control keeps the iterates away from
the boundary, the variance inequality identifies cluster points as fixed
points, and uniqueness finishes the proof.

\begin{example}[The scalar contrast]
In dimension one, $H(x)=c/\sqrt{x}$ and therefore
\[
 K(x)=\frac{c}{\sqrt{x}}x\frac{c}{\sqrt{x}}=c^2=x_\star.
\]
The stabilized update is exact in one step.  The raw map is
$x\mapsto c\sqrt{x}$ and generally needs infinitely many steps.
\end{example}

\subsection{What is known about the raw iteration}

Ruschendorf--Uckelmann considered an equivalent square-root iteration in the
multi-coupling formulation \cite{RuschendorfUckelmann2002}.  Alvarez-Esteban
et al. explicitly distinguish it from \eqref{eq:K-update}, label it their
iteration (27), and state that no theoretical consistency result was available
for it \cite[Section 5]{AlvarezEsteban2016}.  They report satisfactory
two-dimensional examples and quote unfavorable behavior for some
three-dimensional initializations.  These are numerical observations, not a
rigorous divergent orbit or cycle.

As of the literature audit for this note, no primary theorem was located that
proves a Banach contraction of the raw map $\Psi$, or of a fixed iterate
$\Psi^k$, on every invariant spectral compact.  Later global and linear-rate
results concern $K$ or damped Riemannian gradient descent, not raw Picard:
\begin{itemize}
  \item Chewi--Maunu--Rigollet--Stromme prove global rates for Bures gradient
  descent and stochastic gradient descent using a Polyak--Lojasiewicz
  inequality \cite{ChewiEtAl2020}.
  \item Altschuler--Chewi--Gerber--Stromme obtain dimension-free rates under
  uniform eigenvalue bounds.  For a damped step and condition number
  $\kappa$, the objective gap decays like
  $\exp(-cT/\kappa^{5/2})$ with a universal constant $c>0$
  \cite{AltschulerEtAl2021}.
  \item Zemel--Panaretos connect the update with Procrustes alignment and
  Wasserstein gradient descent \cite{ZemelPanaretos2019}.
  \item Brahmachari--Rubboli--Tomamichel prove exponential convergence for a
  broader Bures projection fixed-point algorithm that specializes to the
  Alvarez-Esteban $K$ update \cite{BrahmachariEtAl2025}.
\end{itemize}

\subsection{Other computational formulations}

The Bures fidelity has the semidefinite representation
\begin{equation}\label{eq:fidelity-sdp}
 \tr\bigl(X^{1/2}\Sigma_sX^{1/2}\bigr)^{1/2}
 =\max_{Z_s}\left\{\tr Z_s:
 \begin{pmatrix}X&Z_s\\Z_s^\top&\Sigma_s\end{pmatrix}\succeq0\right\}.
\end{equation}
Consequently, \eqref{eq:barycenter-objective} is an SDP:
\begin{equation}\label{eq:barycenter-sdp}
 \begin{aligned}
 \min_{X,Z_1,\ldots,Z_N}\quad&
 \sum_s\lambda_s\bigl(\tr X+\tr\Sigma_s-2\tr Z_s\bigr),\\
 \text{subject to}\quad&
 \begin{pmatrix}X&Z_s\\Z_s^\top&\Sigma_s\end{pmatrix}\succeq0,
 \qquad s=1,\ldots,N.
 \end{aligned}
\end{equation}
It is robust but expensive at large $d$ or $N$.  Other choices include
Euclidean projected gradient, Bures Riemannian gradient or trust-region
methods, stochastic gradient for population barycenters, and generalized
Procrustes alternating alignment.  For the Gaussian covariance problem, the
stabilized update \eqref{eq:K-update} remains especially simple: each step
requires matrix square roots and linear solves but no step-size selection.

\section{What a contraction statement must say}

\begin{definition}
Let $(D,d)$ be a metric space.  A map $G:D\to D$ is a strict contraction if
there is $q<1$ such that
\[
 d(G(X),G(Y))\leq qd(X,Y),\qquad X,Y\in D.
\]
An iterate $G^k$ is an eventual contraction if this holds for some finite $k$.
\end{definition}

Three quantifiers are distinct:
\begin{enumerate}
  \item for each fixed data set and each fixed invariant compact $D$, there
  exists a $k=k(D,\{\Sigma_s\})$;
  \item for each fixed data set, one $k$ works for every invariant compact;
  \item one universal $k$ works for every data set and compact.
\end{enumerate}
The scalar calculation below proves the first and disproves the second and
third in their unrestricted form.

Metric choice matters.  The Frobenius, operator, affine-invariant,
log-Euclidean, Thompson, and Bures metrics can give different Lipschitz
constants.  Equivalent norms on a finite-dimensional compact have the same
topology, but ``constant $<1$'' is not invariant under a change of norm.

\begin{proposition}[Why eventual contraction would settle raw Picard]
Let $D$ be complete and forward invariant under a continuous map $G$.  If
$G^k:D\to D$ is a contraction for some $k$, then $G$ has a unique fixed point
$p$ in $D$, and every raw orbit $G^n(X)$ converges to $p$.
\end{proposition}

\begin{proof}
Banach's theorem gives a unique fixed point $p$ of $G^k$ and convergence of
$G^{kn}(X)$ to $p$.  Since $G(p)$ is also fixed by $G^k$,
\[
 G^k(G(p))=G(G^k(p))=G(p),
\]
uniqueness implies $G(p)=p$.  Continuity then gives
$G^{kn+r}(X)\to G^r(p)=p$ for each $r=0,\ldots,k-1$.
\end{proof}

Thus a rigorous nonconvergent raw orbit inside an invariant compact would
disprove eventual contraction there.  The historical numerical remarks do not
yet provide such a mathematical counterexample.

\section{Exact scalar, isotropic, and commuting analysis}

\subsection{Dimension one}

Write the scalar covariances as $\sigma_s>0$ and define
\begin{equation}\label{eq:scalar-c}
 c:=\sum_s\lambda_s\sqrt{\sigma_s},\qquad x_\star:=c^2.
\end{equation}
Then $\Psi(x)=c\sqrt{x}$.

\begin{theorem}[Exact iterates and exact Euclidean constant]\label{thm:scalar-exact}
For every $k\geq1$,
\begin{equation}\label{eq:scalar-iterate}
 \Psi^k(x)=x_\star^{1-2^{-k}}x^{2^{-k}},
\end{equation}
and on $[m,M]\subset(0,\infty)$,
\begin{equation}\label{eq:scalar-qk}
 \Lip\bigl(\Psi^k;[m,M]\bigr)
 =q_k:=2^{-k}
 \left(\frac{x_\star}{m}\right)^{1-2^{-k}}.
\end{equation}
The interval is forward invariant exactly when
$m\leq x_\star\leq M$.  On every such fixed interval, $q_k\to0$, so
$\Psi^k$ is a contraction for all sufficiently large $k$.
\end{theorem}

\begin{proof}
The identity is true for $k=1$.  If it holds with exponent $p=2^{-k}$, then
\[
 \Psi^{k+1}(x)
 =\sqrt{x_\star}\,
 \bigl(x_\star^{1-p}x^p\bigr)^{1/2}
 =x_\star^{1-p/2}x^{p/2}.
\]
This proves \eqref{eq:scalar-iterate}.  Its derivative is
\[
 (\Psi^k)'(x)=2^{-k}x_\star^{1-2^{-k}}
 x^{2^{-k}-1}.
\]
It is positive and decreasing, so its maximum on $[m,M]$ is attained at $m$,
which gives \eqref{eq:scalar-qk}.  Since $\Psi$ is increasing, invariance is
equivalent to $\Psi(m)\geq m$ and $\Psi(M)\leq M$, or
$m\leq c^2\leq M$.  Finally $2^{-k}\to0$ while
$(x_\star/m)^{1-2^{-k}}\to x_\star/m$, hence $q_k\to0$.
\end{proof}

\begin{corollary}[No uniform iterate]\label{cor:no-uniform-k}
For every prescribed $k$, there is an invariant compact interval on which
$\Psi^k$ is not a Euclidean contraction.
\end{corollary}

\begin{proof}
Fix $x_\star$ and $M\geq x_\star$.  Formula
\eqref{eq:scalar-qk} exceeds one when $m/x_\star$ is sufficiently small, while
$m\leq x_\star$ keeps the interval invariant.
\end{proof}

For $k=1$, the constant is
\begin{equation}\label{eq:q1-scalar}
 q_1=\frac12\sqrt{\frac{x_\star}{m}},
\end{equation}
so raw Picard expands near the lower end whenever $m<x_\star/4$.

In contrast, the scalar Thompson/log metric gives a global exact contraction:
\begin{equation}\label{eq:scalar-log-contraction}
 \abs{\log\Psi(x)-\log\Psi(y)}
 =\frac12\abs{\log x-\log y}.
\end{equation}
This illustrates why the metric cannot be omitted from the claim.

\subsection{Isotropic data in arbitrary dimension}

Suppose $\Sigma_s=\sigma_s\Id$.  Then
\begin{equation}\label{eq:isotropic-Psi}
 \Psi(X)=cX^{1/2},\qquad
 \Psi^k(X)=x_\star^{1-2^{-k}}X^{2^{-k}}.
\end{equation}

\begin{proposition}[Exact isotropic compact-box constant]
On $\cK_{m,M}$,
\begin{equation}\label{eq:isotropic-q}
 \Lip_F(\Psi^k;\cK_{m,M})
 =2^{-k}\left(\frac{x_\star}{m}\right)^{1-2^{-k}}.
\end{equation}
Thus the scalar conclusions, including nonexistence of a universal $k$, hold
in every matrix dimension.
\end{proposition}

\begin{proof}
For $0<p<1$, the Fr\'echet derivative of $X\mapsto X^p$ in an eigenbasis has
coefficients equal to divided differences of $t^p$.  On $[m,M]$ their maximum
is $pm^{p-1}$ by concavity.  This gives the upper bound in
\eqref{eq:isotropic-q}; the scalar direction at $X=m\Id$ attains it.
\end{proof}

\subsection{Commuting inputs and commuting iterates}

Suppose all inputs and the initial matrix commute.  In a common orthogonal
basis, write
\[
 \Sigma_s=\diag(\sigma_{s,1},\ldots,\sigma_{s,d}),\qquad
 X=\diag(x_1,\ldots,x_d),
\]
and
\begin{equation}\label{eq:commuting-cj}
 c_j:=\sum_s\lambda_s\sqrt{\sigma_{s,j}},
 \qquad x_{\star,j}=c_j^2.
\end{equation}
Then the coordinates decouple:
\begin{equation}\label{eq:commuting-iterate}
 (\Psi^k(X))_{jj}=x_{\star,j}^{1-2^{-k}}x_j^{2^{-k}}.
\end{equation}
On a diagonal rectangle $m_j\leq x_j\leq M_j$, the exact Frobenius
Lipschitz constant is
\begin{equation}\label{eq:commuting-q}
 \max_j 2^{-k}
 \left(\frac{x_{\star,j}}{m_j}\right)^{1-2^{-k}}.
\end{equation}
It tends to zero on every fixed rectangle bounded away from zero.  In
coordinatewise log distance the one-step factor is exactly $1/2$.

The preceding argument restricts both points to the commuting algebra.  The
next result studies arbitrary symmetric perturbations around the commuting
barycenter and is stronger.

\begin{theorem}[Local derivative for commuting data]\label{thm:commuting-local}
Assume the $\Sigma_s$ are diagonal in a common basis.  At the barycenter
$X_\star=\diag(c_1^2,\ldots,c_d^2)$, the derivative
$J=D\Psi(X_\star)$ is diagonal on the standard symmetric matrix basis.
For a diagonal direction $E_{ii}$ its eigenvalue is $1/2$.  For an off-diagonal
direction $E_{ij}+E_{ji}$, let
$u_s=\sqrt{\sigma_{s,i}}$ and $v_s=\sqrt{\sigma_{s,j}}$.  The eigenvalue is
\begin{equation}\label{eq:theta-ij}
 \theta_{ij}=
 \frac{1}{c_i+c_j}\sum_s\lambda_s
 \frac{c_i u_s^2+c_jv_s^2}{c_iu_s+c_jv_s}.
\end{equation}
It satisfies
\begin{equation}\label{eq:theta-bounds}
 \frac12\leq\theta_{ij}<1.
\end{equation}
Hence $\norm{D\Psi(X_\star)}_{F\to F}<1$, and $\Psi$ is a Frobenius
contraction on some full, not merely diagonal, neighborhood of $X_\star$.
\end{theorem}

\begin{proof}
The derivative formula established in \cref{thm:frechet-derivative} below
preserves each coordinate matrix when all matrices at the base point are
diagonal.  Substitution of
$X_\star^{1/2}=\diag(c_1,\ldots,c_d)$ gives
\eqref{eq:theta-ij}; the diagonal coefficient reduces to
$\sum_s\lambda_su_s/(2c_i)=1/2$.

For the upper bound, pointwise in $s$,
\[
 \frac{c_i u_s^2+c_jv_s^2}{c_iu_s+c_jv_s}
 <u_s+v_s,
\]
because the difference after multiplication by the positive denominator is
$(c_i+c_j)u_sv_s>0$.  Averaging gives a strict upper bound $c_i+c_j$, and
division by $c_i+c_j$ gives $\theta_{ij}<1$.

For the lower bound, weighted Cauchy--Schwarz gives
\[
 c_i u_s^2+c_jv_s^2\geq
 \frac{(c_iu_s+c_jv_s)^2}{c_i+c_j}.
\]
Therefore
\[
 \theta_{ij}\geq
 \frac{c_i^2+c_j^2}{(c_i+c_j)^2}\geq\frac12.
\]
The derivative is self-adjoint and diagonal in an orthonormal Frobenius basis,
so its operator norm is its largest eigenvalue, strictly below one.  Continuity
of $D\Psi$ yields a neighborhood on which the derivative norm remains below
one; the mean-value formula then gives a contraction.
\end{proof}

\begin{example}[Slow off-diagonal modes]
For a single input $A=\diag(a_1,\ldots,a_d)$, the fixed point is $A$ and
\begin{equation}\label{eq:single-input-theta}
 D\Psi(A)[E_{ij}+E_{ji}]
 =\left(1-\frac{\sqrt{a_ia_j}}{a_i+a_j}\right)(E_{ij}+E_{ji}).
\end{equation}
For $i\neq j$, the coefficient equals $1/2$ when $a_i=a_j$ and approaches
$1$ as $a_i/a_j\to\infty$; diagonal modes separately have eigenvalue $1/2$.
Thus even local raw Picard convergence can be arbitrarily
slow with ill-conditioning; no condition-number-free strict factor is
possible.
\end{example}

\section{Full noncommutative differential calculus}

\subsection{Derivative through Sylvester equations}

For $Z\succ0$, define the Lyapunov/Sylvester operator
\begin{equation}\label{eq:lyapunov}
 \cL_Z(U):=ZU+UZ.
\end{equation}
It is invertible on symmetric matrices.  Differentiating $S^2=X$ shows that
the Fr\'echet derivative of the matrix square root is
\begin{equation}\label{eq:dsqrt}
 D(X^{1/2})[H]=\cL_{X^{1/2}}^{-1}(H).
\end{equation}

\begin{theorem}[Exact derivative of the raw map]\label{thm:frechet-derivative}
Fix $X\succ0$ and a symmetric perturbation $H$.  Put
\[
 S=X^{1/2},\qquad R_s=(S\Sigma_sS)^{1/2}.
\]
Let $U$ and $V_s$ be the unique symmetric solutions of
\begin{align}
 SU+US&=H,\label{eq:sylvester-U}\\
 R_sV_s+V_sR_s&=U\Sigma_sS+S\Sigma_sU.
 \label{eq:sylvester-V}
\end{align}
Then
\begin{equation}\label{eq:Dpsi}
 \boxed{D\Psi(X)[H]=\sum_s\lambda_sV_s.}
\end{equation}
\end{theorem}

\begin{proof}
Equation \eqref{eq:sylvester-U} is the derivative of $S^2=X$.  For
$Y_s=S\Sigma_sS$, the product rule gives
\[
 DY_s[H]=U\Sigma_sS+S\Sigma_sU.
\]
Differentiating $R_s^2=Y_s$ gives \eqref{eq:sylvester-V}.  Sum with the fixed
weights.
\end{proof}

The formula is numerically stable: diagonalize the two positive matrices or
solve the Sylvester equations directly.  It avoids finite differences and
allows exact evaluation of the Jacobian on the $d(d+1)/2$-dimensional
symmetric space.

\subsection{A conservative compact-box Lipschitz bound}

\begin{lemma}[Square-root derivative norm]\label{lem:sqrt-deriv-bound}
For $Z\succ0$,
\begin{equation}\label{eq:lyapunov-inverse-bound}
 \norm{\cL_Z^{-1}}_{F\to F}=\frac{1}{2\lambda_{\min}(Z)}.
\end{equation}
\end{lemma}

\begin{proof}
In an eigenbasis of $Z=\diag(z_1,\ldots,z_d)$,
$(\cL_Z^{-1}H)_{ij}=H_{ij}/(z_i+z_j)$.  The largest coefficient is
$1/(2\min_i z_i)$.
\end{proof}

\begin{proposition}[One-step analytic bound]\label{prop:one-step-bound}
Under \eqref{eq:data-bounds}, on $\cK_{m,M}$,
\begin{equation}\label{eq:LmM}
 \Lip_F(\Psi;\cK_{m,M})\leq
 L(m,M):=\frac{\sqrt M}{2m}
 \sum_s\lambda_s\frac{b_s}{\sqrt{a_s}}.
\end{equation}
\end{proposition}

\begin{proof}
Let $\norm{H}_F=1$.  From \eqref{eq:sylvester-U},
\[
 \norm U_F\leq\frac{1}{2\sqrt m}.
\]
Also $\norm S_{\mathrm{op}}\leq\sqrt M$ and
$\norm{\Sigma_s}_{\mathrm{op}}\leq b_s$, hence
\[
 \norm{U\Sigma_sS+S\Sigma_sU}_F
 \leq b_s\sqrt{\frac{M}{m}}.
\]
Since $S\Sigma_sS\succeq a_sm\Id$, we have
$R_s\succeq\sqrt{a_sm}\Id$.  A second use of
\cref{lem:sqrt-deriv-bound} gives
\[
 \norm{V_s}_F\leq
 \frac{b_s\sqrt M}{2m\sqrt{a_s}}.
\]
After summing, this bounds $\norm{D\Psi(X)}_{F\to F}$ uniformly on the convex
box.  Integrating the derivative along the segment from $X$ to $Y$ proves
\eqref{eq:LmM}.
\end{proof}

This bound is deliberately elementary and can be loose.  Its value is that it
is explicit and composable.

\subsection{A rigorous sufficient theorem for an iterate}

Starting from $m_0=m$, $M_0=M$, define
\begin{equation}\label{eq:box-recursion}
 m_{j+1}=c_-\sqrt{m_j},\qquad
 M_{j+1}=c_+\sqrt{M_j}.
\end{equation}
Equivalently,
\begin{equation}\label{eq:box-explicit}
 m_j=c_-^{2(1-2^{-j})}m^{2^{-j}},\qquad
 M_j=c_+^{2(1-2^{-j})}M^{2^{-j}}.
\end{equation}
By \cref{lem:spectral-enclosure},
$\Psi^j(\cK_{m,M})\subset\cK_{m_j,M_j}$.

\begin{theorem}[Compact-box contraction certificate]\label{thm:compact-certificate}
Let
\begin{equation}\label{eq:Qk}
 Q_k:=\prod_{j=0}^{k-1}L(m_j,M_j),
\end{equation}
where $L$ is defined in \eqref{eq:LmM}.  Then
\begin{equation}\label{eq:iterate-Lip-bound}
 \Lip_F(\Psi^k;\cK_{m,M})\leq Q_k.
\end{equation}
In particular, $Q_k<1$ is a rigorous sufficient certificate that $\Psi^k$ is
a Frobenius contraction on the box.  If the box is invariant by
\eqref{eq:invariance-condition}, Banach's theorem applies.

Moreover, put
\begin{equation}\label{eq:L-infty}
 L_\infty:=
 \frac{c_+}{2c_-^2}
 \sum_s\lambda_s\frac{b_s}{\sqrt{a_s}}.
\end{equation}
If $L_\infty<1$, then $Q_k\to0$ and some iterate is a contraction on every
fixed box $\cK_{m,M}$.
\end{theorem}

\begin{proof}
Each factor in the derivative chain for $\Psi^k$ is evaluated in the
corresponding image box.  Apply \cref{prop:one-step-bound} successively to
obtain \eqref{eq:iterate-Lip-bound}.  From \eqref{eq:box-explicit},
$m_j\to c_-^2$ and $M_j\to c_+^2$, so
$L(m_j,M_j)\to L_\infty$.  If the limit is below one, all sufficiently late
factors are bounded by some $q<1$, and the finite product of early factors
times $q^{k-j_0}$ tends to zero.
\end{proof}

The condition $L_\infty<1$ is sufficient, not necessary.  Failure of this
test says only that the elementary eigenvalue estimates lost too much
information.  The exact derivative product may still contract.

\subsection{The sharp local criterion}

Let $X_\star$ be the barycenter and $J=D\Psi(X_\star)$.  Homogeneity gives a
universal eigenpair:
\begin{equation}\label{eq:radial-eigenpair}
 J[X_\star]=\frac12X_\star.
\end{equation}
Indeed, differentiate $\Psi(tX_\star)=\sqrt t X_\star$ at $t=1$.

\begin{theorem}[Local eventual contraction criterion]\label{thm:local-criterion}
Let $G$ be a $C^1$ map on a finite-dimensional normed space with fixed point
$p$, and let $J=DG(p)$.
\begin{enumerate}
  \item If $\norm J<1$ in the chosen induced norm, then $G$ is a contraction
  on a sufficiently small invariant ball around $p$.
  \item If the spectral radius $\rhoop(J)<1$, then for every fixed induced
  norm there is a $k$ such that $G^k$ is a contraction on a sufficiently small
  invariant ball around $p$.
  \item If $\rhoop(J)\geq1$, no iterate can be a strict local contraction in
  any norm.
\end{enumerate}
Applied to $G=\Psi$, the best possible local factor is at least $1/2$ because
of \eqref{eq:radial-eigenpair}.
\end{theorem}

\begin{proof}
The first statement follows from continuity of $DG$ and the mean-value
formula.  If $\rhoop(J)<1$, then $J^k\to0$, hence $\norm{J^k}<1$ for some
$k$; use the first statement for $G^k$.  Conversely, a local contraction of
$G^k$ implies $\norm{J^k}<1$, so
$\rhoop(J)^k=\rhoop(J^k)<1$.
\end{proof}

\subsection{What remains open for the raw map}

The general noncommuting problem can now be stated without ambiguity.

\begin{quote}
For fixed positive definite data and a fixed forward-invariant spectral box,
is there always a finite $k$ for which
\[
 \sup_{X\in\cK_{m,M}}
 \norm{D\Psi(\Psi^{k-1}X)\cdots D\Psi(X)}_{F\to F}<1?
\]
\end{quote}

The exact scalar and commuting results support a positive answer in those
regimes.  The compact-box theorem proves it under an explicit conditioning
condition.  The literature located for this study does not settle the fully
noncommuting statement.  A useful first local target is
\begin{equation}\label{eq:local-open-target}
 \rhoop(D\Psi(X_\star))<1
 \quad\text{for arbitrary noncommuting positive definite inputs}.
\end{equation}
Even a proof of \eqref{eq:local-open-target} would be local; extending it
uniformly over a whole spectral box would require control of derivative
products along every raw orbit.

\section{Practical derivative test and reference algorithms}

For a concrete data set, the following procedure separates provable facts
from numerical evidence.

\paragraph{Step 1: compute the barycenter reliably.}
Use the stabilized update
\[
 X\leftarrow X^{-1/2}\Psi(X)^2X^{-1/2}
\]
until either the fixed residual $\norm{\Psi(X)-X}_F/\norm X_F$ or the Bures
step is below tolerance.  Symmetrize after floating-point products and compute
square roots by an eigendecomposition with positive eigenvalue checks.

\paragraph{Step 2: assemble the exact Jacobian.}
Choose the orthonormal symmetric basis
\[
 E_{ii},\qquad (E_{ij}+E_{ji})/\sqrt2\quad(i<j).
\]
For each basis element solve \eqref{eq:sylvester-U} and
\eqref{eq:sylvester-V}; project the result onto the basis.  This forms the
$d(d+1)/2$ square matrix representing $D\Psi(X_\star)$.

\paragraph{Step 3: distinguish three diagnostics.}
\begin{itemize}
  \item $\rhoop(J)<1$ certifies local eventual contraction.
  \item $\norm J_2<1$ certifies local one-step Frobenius contraction.
  \item Neither diagnostic certifies contraction on a large compact box.
  For that, bound the derivative chain along the full reachable set, using
  \cref{thm:compact-certificate} or sharper interval/spectral estimates.
\end{itemize}

\paragraph{Step 4: search for a rigorous obstruction if needed.}
A pair $X,Y$ with
\[
 \norm{\Psi^k(X)-\Psi^k(Y)}_F\geq\norm{X-Y}_F
\]
disproves $k$-step Frobenius contraction on any domain containing them, but
does not disprove eventual contraction for a larger iterate.  A verified
nontrivial periodic orbit inside an invariant compact would disprove every
eventual-contraction theorem there.

\section{Entropic and Sinkhorn Gaussian barycenters}

\subsection{Three objectives that must not be conflated}

Fix the convention
\begin{equation}\label{eq:EOT-convention}
 \mathrm{OT}_\eps(\mu,\nu)
 :=\inf_{\pi\in\Pi(\mu,\nu)}
 \left\{\frac12\int\norm{x-y}_2^2\dd\pi
 +\eps\KL(\pi\mid\mu\otimes\nu)\right\}.
\end{equation}
There are at least three common barycenter objectives.

\begin{enumerate}
  \item The \emph{product-reference raw entropic barycenter} minimizes
  $\sum_s\lambda_s\mathrm{OT}_\eps(\mu,\nu_s)$.  It has a shrinkage bias and
  can collapse.
  \item The \emph{Lebesgue/Brownian-reference Schrodinger barycenter} differs
  by a marginal entropy term.  With
  $H(\mu)=\int\rho\log\rho$, the family
  \[
    \sum_s\lambda_s\mathrm{OT}_\eps(\mu,\nu_s)+\tau H(\mu)
  \]
  has product reference at $\tau=0$ and the Lebesgue/Brownian convention at
  $\tau=\eps$.  The latter has a blurring bias.
  \item The \emph{debiased Sinkhorn divergence}
  \begin{equation}\label{eq:sinkhorn-divergence}
    S_\eps(\mu,\nu)=\mathrm{OT}_\eps(\mu,\nu)
    -\frac12\mathrm{OT}_\eps(\mu,\mu)
    -\frac12\mathrm{OT}_\eps(\nu,\nu)
  \end{equation}
  cancels the reference-dependent marginal entropy terms and much of the
  entropic bias.
\end{enumerate}

This taxonomy and the doubly regularized interpolation are developed in
\cite{Chizat2025}.  Pairwise Gaussian entropic OT has a Gaussian optimal
coupling, but that pairwise statement alone does not prove that a barycenter
over all probability measures is Gaussian.

\subsection{Gaussian covariance equations}

Let the inputs be $\cN(m_s,A_s)$ and write $\delta=\eps/2$.  The mean is again
$m=\sum_s\lambda_sm_s$.  Within the Gaussian manifold, the product-reference
raw covariance equation is
\begin{equation}\label{eq:entropic-product}
 B=\sum_s\lambda_s
 \bigl(B^{1/2}A_sB^{1/2}+\delta^2\Id\bigr)^{1/2}-\delta\Id.
\end{equation}
More generally, adding $\tau H(\mu)$ gives the Gaussian stationarity equation
\begin{equation}\label{eq:doubly-regularized}
 B=\sum_s\lambda_s
 \bigl(B^{1/2}A_sB^{1/2}+\delta^2\Id\bigr)^{1/2}
 +(\tau-\delta)\Id.
\end{equation}
Thus the Lebesgue/Brownian choice $\tau=\eps$ changes $-\delta\Id$ into
$+\delta\Id$.  The product-reference equation and numerical fixed-point
iteration are derived by Mallasto--Gerolin--Minh
\cite{MallastoGerolinMinh2022}; their general multivariate barycenter is
Gaussian-restricted, and they leave convergence of the covariance iteration
open.

For the debiased Sinkhorn divergence, Janati--Muzellec--Peyre--Cuturi prove a
stronger arbitrary-dimensional result: minimization over a class of
sub-Gaussian probability measures has a Gaussian minimizer.  Its covariance
satisfies
\begin{equation}\label{eq:sinkhorn-covariance}
 \sum_s\lambda_s
 \bigl(B^{1/2}A_sB^{1/2}+\delta^2\Id\bigr)^{1/2}
 =\bigl(B^2+\delta^2\Id\bigr)^{1/2}.
\end{equation}
See \cite[Theorem 2]{JanatiEtAl2020UnbalancedGaussian}.  An associated
covariance Picard map is
\begin{equation}\label{eq:sinkhorn-map}
 \cT_\eps(B):=
 \left[
 \left(\sum_s\lambda_s
 (B^{1/2}A_sB^{1/2}+\delta^2\Id)^{1/2}\right)^2
 -\delta^2\Id
 \right]^{1/2}.
\end{equation}
The existence proof uses an invariant spectral set and Brouwer, not a
contraction of \eqref{eq:sinkhorn-map}.  Uniqueness within broad Gaussian
covariance classes, including infinite-dimensional versions, is studied in
\cite{Minh2023}.

\subsection{Scalar anti-contraction and eventual-contraction checks}

For scalar variances $a_s,b>0$, define
\begin{align}
 F_P(b)&=\sum_s\lambda_s\sqrt{a_sb+\delta^2}-\delta,
 \label{eq:scalar-entropic-P}\\
 F_L(b)&=\sum_s\lambda_s\sqrt{a_sb+\delta^2}+\delta,
 \label{eq:scalar-entropic-L}\\
 F_S(b)&=\sqrt{\left(\sum_s\lambda_s
 \sqrt{a_sb+\delta^2}\right)^2-\delta^2}.
 \label{eq:scalar-entropic-S}
\end{align}
The first two derivatives are
\begin{equation}\label{eq:entropic-derivative}
 F_P'(b)=F_L'(b)=
 \sum_s\lambda_s\frac{a_s}{2\sqrt{a_sb+\delta^2}},
\end{equation}
which can exceed one near $b=0$ for legitimate parameters.  Thus entropic
regularization does not automatically produce a Euclidean contraction.

For equal inputs $a_s=a$, the debiased map simplifies exactly to
\begin{equation}\label{eq:equal-sinkhorn-scalar}
 F_S(b)=\sqrt{ab},
\end{equation}
independently of $\eps$.  It therefore reproduces all the scalar raw-map
phenomena from \cref{thm:scalar-exact}: non-Lipschitz behavior at the boundary,
possible one-step expansion on a wide interval, and eventual contraction on
each fixed invariant interval bounded away from zero.  No contraction theorem
uniform in the lower spectral bound can hold.

For equal isotropic inputs $a\Id$, the doubly regularized covariance is
explicit \cite{Chizat2025}:
\begin{equation}\label{eq:isotropic-double-reg}
 b=\frac{\left(a+\sqrt{(a-\eps)^2+4a\tau}\right)^2-\eps^2}{4a}.
\end{equation}
Hence product reference gives $b=(a-\eps)_+$, while the
Lebesgue/Brownian choice gives $b=a+\eps$.  The two common raw conventions
have opposite biases.

\subsection{Algorithms and theorem scopes}

\begin{itemize}
  \item On a fixed grid, raw entropic barycenters can be solved by matrix
  scaling/iterative Bregman projections; convergence follows from the KL
  projection framework \cite{CuturiDoucet2014,BenamouEtAl2015}.
  \item Debiased fixed-grid Sinkhorn barycenters have efficient scaling
  algorithms, but empirical linear behavior should not be restated as a
  proof of contraction of \eqref{eq:sinkhorn-map}
  \cite{JanatiCuturiGramfort2020}.
  \item Free-support alternatives include Frank--Wolfe and functional
  gradient schemes with sublinear objective or stationarity guarantees
  \cite{LuiseEtAl2019,ShenEtAl2020}.
  \item A contraction proved for an inner pairwise Gaussian dual-potential
  iteration is not a contraction of the barycenter covariance maps
  \eqref{eq:entropic-product}, \eqref{eq:doubly-regularized}, or
  \eqref{eq:sinkhorn-map}.
\end{itemize}

The safest preservation statements are therefore:
\begin{itemize}
  \item product-reference raw, general $d$: Gaussian-restricted equation;
  \item Lebesgue raw: full Gaussian closure is known in particular for two
  multivariate inputs and for multiple inputs in one dimension
  \cite{CumingsMenonShin2020,JanatiCuturiGramfort2020};
  \item debiased Sinkhorn: a full sub-Gaussian closure theorem in arbitrary
  dimension \cite{JanatiEtAl2020UnbalancedGaussian}.
\end{itemize}

\section{Unbalanced Gaussian barycenters and WFR/HK}

\subsection{Ambient Hellinger--Kantorovich geometry}

The dynamic Wasserstein--Fisher--Rao, also called Hellinger--Kantorovich,
action combines transport and reaction.  One normalization is
\begin{equation}\label{eq:wfr-dynamic}
 \partial_t\rho_t+\nabla\!\cdot(\rho_tv_t)=r_t\rho_t,
 \qquad
 \int_0^1\int\left(\frac1\alpha\norm{v_t}^2
 +\frac1\beta r_t^2\right)\rho_t\dd x\dd t,
\end{equation}
with $\alpha,\beta>0$.  Equivalent static entropy-transport formulations use
a truncated cosine cost; see
\cite{LieroMielkeSavare2018,ChizatEtAl2018}.

The term ``Gaussian Hellinger--Kantorovich'' is potentially misleading.  In
one usage, GHK refers to a quadratic-cost KL-unbalanced transport discrepancy;
``Gaussian'' describes the ground kernel, not closure of Gaussian measures.
Liero--Mielke--Savare show that, for
$g(r)=\arccos(e^{-r^2/2})$, this discrepancy is related by
$\mathrm{GHK}_d=\mathrm{HK}_{g\circ d}$; ordinary HK is the induced length
metric in the appropriate setting \cite[Section 7.8]{LieroMielkeSavare2018}.

\subsection{An exact negative result at the Hellinger endpoint}

At the pure-reaction endpoint, the square-root embedding linearizes the
squared Hellinger distance.  For nonnegative densities $\rho_s$, the
unconstrained-mass barycenter is
\begin{equation}\label{eq:hellinger-barycenter}
 \boxed{\rho_H^\star(x)=
 \left(\sum_s\lambda_s\sqrt{\rho_s(x)}\right)^2.}
\end{equation}

\begin{proposition}[Gaussian inputs are not preserved by Hellinger barycenters]
Let $\rho_s=a_sg_s$, where $a_s>0$ and
$g_s$ is the density of $\cN(m_s,\Sigma_s)$.  Then
\begin{equation}\label{eq:hellinger-mixture}
 \rho_H^\star
 =\sum_{s,t}\lambda_s\lambda_t\sqrt{a_sa_t}\sqrt{g_sg_t}.
\end{equation}
Each $\sqrt{g_sg_t}$ is a positive constant times a Gaussian density with
\begin{align}
 \overline\Sigma_{st}&=2(\Sigma_s^{-1}+\Sigma_t^{-1})^{-1},
 \label{eq:cross-cov}\\
 \overline m_{st}&=(\Sigma_s^{-1}+\Sigma_t^{-1})^{-1}
 (\Sigma_s^{-1}m_s+\Sigma_t^{-1}m_t).
 \label{eq:cross-mean}
\end{align}
Thus \eqref{eq:hellinger-mixture} is generically a finite Gaussian mixture,
not a single Gaussian.
\end{proposition}

\begin{proof}
Minimizing
$\sum_s\lambda_s\norm{q-\sqrt{\rho_s}}_{L^2}^2$ over $q\geq0$ gives
$q=\sum_s\lambda_s\sqrt{\rho_s}$, proving
\eqref{eq:hellinger-barycenter}.  Expanding the square gives
\eqref{eq:hellinger-mixture}.  Completing the square in the product of the two
Gaussian half-densities gives \eqref{eq:cross-cov} and
\eqref{eq:cross-mean}.  Unless all components coincide in a degenerate way,
the resulting mixture is not one Gaussian.
\end{proof}

The normalizing constant of the cross term is the Bhattacharyya affinity
\begin{equation}\label{eq:bhattacharyya}
 \int\sqrt{g_sg_t}
 =\frac{\det(\Sigma_s)^{1/4}\det(\Sigma_t)^{1/4}}
 {\det((\Sigma_s+\Sigma_t)/2)^{1/2}}
 \exp\left[-\frac18(m_s-m_t)^\top
 \left(\frac{\Sigma_s+\Sigma_t}{2}\right)^{-1}(m_s-m_t)\right].
\end{equation}

\subsection{Why Gelbrich's Gaussianization proof cannot transfer}

The obstruction is stronger than lack of a known moment formula.

\begin{proposition}[Moment Gaussianization is not HK-continuous]\label{prop:HK-counterexample}
Let $\gamma=\cN(0,1)$ on $\R$ and
\begin{equation}\label{eq:HK-mixture-counterexample}
 \mu_\eps=(1-\eps)\gamma+eps\cN(\eps^{-1/2},1).
\end{equation}
Then the Hellinger and HK distances from $\mu_\eps$ to $\gamma$ tend to zero,
but
\begin{equation}\label{eq:HK-G-counterexample}
 \cG(\mu_\eps)=\cN(\sqrt\eps,2-\eps)
 \longrightarrow\cN(0,2)\neq\gamma.
\end{equation}
All covariances lie in the compact interval $[1,2]$.  Hence moment
Gaussianization is not continuous, and therefore not nonexpansive, for HK even
under a compact covariance bound.
\end{proposition}

\begin{proof}
The $L^1$ distance between $\mu_\eps$ and $\gamma$ is at most $2\eps$.  For
any common dominating measure $\xi$,
\[
 \int\left(\sqrt{\frac{\dd\mu_\eps}{\dd\xi}}
 -\sqrt{\frac{\dd\gamma}{\dd\xi}}\right)^2\dd\xi
 \leq\norm{\mu_\eps-\gamma}_{L^1}\leq2\eps.
\]
HK can realize pure reaction, so it is bounded above, up to normalization, by
the Hellinger distance.  Thus $\mathrm{HK}(\mu_\eps,\gamma)\to0$.

The mean of \eqref{eq:HK-mixture-counterexample} is $\sqrt\eps$, its second
moment is $2$, and its variance is $2-\eps$, proving
\eqref{eq:HK-G-counterexample}.  Since HK is a metric and metrizes the relevant
narrow topology for bounded masses, the distance from
$\cN(\sqrt\eps,2-\eps)$ to $\gamma$ tends to the strictly positive distance
between $\cN(0,2)$ and $\cN(0,1)$.
\end{proof}

The conceptual difference is that $W_2$ convergence includes convergence of
second moments, whereas HK allows a vanishing amount of mass to travel far
away or to be created/destroyed.  Covariance is therefore not controlled by
the ambient HK topology.

\subsection{A formal geodesic obstruction for positive transport and reaction}

For a smooth HK geodesic, an optimal potential satisfies, after a choice of
normalization,
\begin{align}
 \partial_t\rho&=-\alpha\nabla\!\cdot(\rho\nabla\phi)+\beta\rho\phi,
 \label{eq:HK-continuity}\\
 \partial_t\phi+\frac\alpha2\norm{\nabla\phi}^2
 +\frac\beta2\phi^2&=0.
 \label{eq:HK-HJ}
\end{align}

Suppose formally that
$\rho_t=\kappa_t\cN(m_t,\Sigma_t)$ stays Gaussian on a time interval.  Then
$(\partial_t\rho_t)/\rho_t$ is a quadratic polynomial in $x$.  The weighted
elliptic equation \eqref{eq:HK-continuity} has a quadratic finite-energy
solution $\phi_t$ (this can be read from the Gaussian Hermite decomposition).
If the quadratic part is nonzero, the term $\phi_t^2$ in
\eqref{eq:HK-HJ} has a nonzero quartic part, impossible because the other
terms have degree at most two.  Hence the quadratic part vanishes.  If a
linear part remains, $\phi_t^2$ has a nonzero quadratic part while
$\partial_t\phi_t$ is affine and $\norm{\nabla\phi_t}^2$ is constant, again a
contradiction.  Therefore $\phi_t$ is spatially constant: only mass changes,
while mean and covariance stay fixed.

This is a smooth/formal total-geodesy obstruction, not a substitute for a full
weak-geodesic theorem.  It explains structurally why the Wasserstein Gaussian
closure mechanism disappears as soon as the reaction term $\beta\phi^2$ is
present.  A particular HK gradient flow may nevertheless preserve Gaussians;
for example, the HK-Boltzmann flow with a Gaussian target studied in
\cite{LieroMielkeTseZhu2025}.  Invariance of one vector field does not imply
that the Gaussian family is geodesically closed or closed under ambient
barycenters.

\subsection{Three legitimate unbalanced routes}

\paragraph{1. Ambient HK/WFR barycenter.}
Solve in the full measure space.  Existence, multimarginal formulations, and
Dirac structure are studied in
\cite{FrieseckeMatthesSchmitzer2021,ChungPhung2021}.  After discretization,
generalized Sinkhorn/diagonal-scaling algorithms handle entropy-regularized
unbalanced transport and barycenters \cite{ChizatEtAl2018}.  The result is an
unrestricted measure and should not be expected to be Gaussian.

\paragraph{2. Gaussian-restricted BWFR geometry.}
Constrain paths a priori to
\begin{equation}\label{eq:scaled-gaussian-family}
 \rho=\kappa\cN(m,\Sigma).
\end{equation}
The projected/reduced Gaussian structure derived in
\cite{LieroMielkeTseZhu2025} has a co-metric operator which, on a cotangent
triple $(S,u,k)$, takes the form
\begin{equation}\label{eq:gaussian-WFR-cometric}
 K^G_{\alpha,\beta}(\Sigma,m,\kappa)[S,u,k]
 =\left(
 \frac{2\alpha}{\kappa}(S\Sigma+\Sigma S)
 +\frac{2\beta}{\kappa}\Sigma S\Sigma,
 \frac\alpha\kappa u+\frac\beta\kappa\Sigma u,
 \beta\kappa k
 \right).
\end{equation}
The induced path distance $d_G$ is obtained by inverting this positive
operator in the action.  Its mean and mass terms include
\begin{equation}\label{eq:gaussian-WFR-primal-pieces}
 \kappa\dot m^\top(\alpha\Id+\beta\Sigma)^{-1}\dot m
 +\frac{\dot\kappa^2}{\beta\kappa};
\end{equation}
the covariance term is the inverse of the Lyapunov/Riccati operator displayed
in \eqref{eq:gaussian-WFR-cometric}.  Barycenters for $d_G$ are Gaussian by
construction and can be computed by Riemannian shooting, log-map gradient
descent, or direct path optimization.  Since the path class is restricted,
\begin{equation}\label{eq:restricted-upper}
 \mathrm{WFR}_{\mathrm{ambient}}\leq d_G
\end{equation}
for Gaussian endpoints in compatible normalization.  This is a constrained
surrogate, not an ambient preservation theorem.

\paragraph{3. Quadratic-cost KL-unbalanced and semi-unbalanced models.}
Entropic quadratic unbalanced OT between scaled Gaussians has a scaled
Gaussian optimal plan \cite{JanatiEtAl2020UnbalancedGaussian}.  Recent work
also gives closed Gaussian pairwise formulas without coupling entropy for
quadratic-cost KL-unbalanced OT \cite{YangZhang2026}; a Gaussian barycenter
theorem is a separate question.  Semi-unbalanced OT with one hard marginal
does admit a Gaussian barycenter theorem and Bures-gradient algorithms
\cite{NguyenEtAl2026}.  It is asymmetric and is not the ambient HK/WFR metric,
but the hard barycenter marginal restores a moment-projection mechanism and
can be useful for robust averaging.

\section{A research program for the raw map}

The contraction question is now narrow enough to be attacked directly.

\begin{enumerate}
  \item \textbf{Settle the local Jacobian.}  Prove or disprove
  \eqref{eq:local-open-target}.  The Sylvester representation
  \eqref{eq:Dpsi}, the universal radial eigenvalue $1/2$, and the exact
  commuting spectrum \eqref{eq:theta-ij} provide strong structure.

  \item \textbf{Find the right inner product.}  The general Jacobian need not
  be visibly self-adjoint in Frobenius geometry.  A data- or fixed-point-
  dependent Lyapunov inner product might reveal positivity and an upper
  spectral bound.  Any adapted norm must then be related quantitatively to the
  desired metric on the compact domain.

  \item \textbf{Control products, not only factors.}  Individual derivative
  norms may exceed one far from the solution.  Eventual contraction requires
  decay of
  $D\Psi(\Psi^{k-1}X)\cdots D\Psi(X)$ uniformly in $X$.  The scalar proof works
  precisely because early expansion is followed by exponent halving.

  \item \textbf{Exploit spectral trapping.}  The explicit boxes
  \eqref{eq:box-explicit} forget orientation but rapidly stabilize the scale.
  Sharper bounds should propagate eigenvector/alignment information or use
  the transport maps $T_s(X)$ rather than separate min/max eigenvalues.

  \item \textbf{Separate proof goals.}  Global convergence of raw Picard,
  local linear convergence, one-step contraction, and eventual contraction
  on a compact are different statements.  The globally convergent $K$ update
  already solves the computational problem; a raw-$\Psi$ contraction theorem
  would mainly clarify the nonlinear matrix equation and provide an alternate
  solver.
\end{enumerate}

\section{Final perspective}

The Gaussian barycenter covariance equation is well understood as an
optimality condition, and the barycenter itself is unique.  What is not
justified is the inference
\[
 \text{unique fixed point}\quad\Longrightarrow\quad
 \text{raw fixed-point map is a contraction}.
\]
The classical algorithmic literature avoids this gap by iterating the average
optimal transport map, producing $K(X)=X^{-1/2}\Psi(X)^2X^{-1/2}$ and a
variance descent inequality.

For the raw map, the scalar answer is more positive than the initial suspicion:
on each fixed compact interval separated from zero, a high enough iterate is a
strict contraction.  Yet the exact constant proves that no iterate can be
chosen uniformly as the lower bound approaches the PSD boundary.  The same
obstruction survives in every dimension through isotropic data.  In the
commuting case, even arbitrary off-diagonal perturbations are locally
contracted, with an explicit derivative spectrum in $[1/2,1)$.  The exact
Sylvester derivative and compact-box certificate reduce the genuinely open
part to noncommutative uniform control of derivative products.

Entropic regularization does not erase this issue; different entropy
conventions produce different biases and covariance equations, and scalar
maps still need not contract in one step.  In the unbalanced setting the
situation changes more radically: ambient HK/WFR does not control moments,
moment Gaussianization is not continuous, and Hellinger barycenters of
Gaussians are generically Gaussian mixtures.  Gaussian-restricted BWFR or
semi-unbalanced Gaussian models remain useful, provided their restricted or
asymmetric nature is stated explicitly.

\clearpage
\small
\setlength{\parskip}{0.15em}
\begin{thebibliography}{99}

\bibitem{Gelbrich1990}
M. Gelbrich,
``On a Formula for the $L^2$ Wasserstein Metric between Measures on Euclidean
and Hilbert Spaces,''
\emph{Mathematische Nachrichten} 147 (1990), 185--203.
\url{https://doi.org/10.1002/mana.19901470121}

\bibitem{DowsonLandau1982}
D. C. Dowson and B. V. Landau,
``The Fr\'echet Distance between Multivariate Normal Distributions,''
\emph{Journal of Multivariate Analysis} 12(3) (1982), 450--455.
\url{https://doi.org/10.1016/0047-259X(82)90077-X}

\bibitem{OlkinPukelsheim1982}
I. Olkin and F. Pukelsheim,
``The Distance between Two Random Vectors with Given Dispersion Matrices,''
\emph{Linear Algebra and its Applications} 48 (1982), 257--263.
\url{https://doi.org/10.1016/0024-3795(82)90112-4}

\bibitem{GivensShortt1984}
C. R. Givens and R. M. Shortt,
``A Class of Wasserstein Metrics for Probability Distributions,''
\emph{Michigan Mathematical Journal} 31(2) (1984), 231--240.
\url{https://doi.org/10.1307/mmj/1029003026}

\bibitem{AguehCarlier2011}
M. Agueh and G. Carlier,
``Barycenters in the Wasserstein Space,''
\emph{SIAM Journal on Mathematical Analysis} 43(2) (2011), 904--924.
\url{https://doi.org/10.1137/100805741}

\bibitem{KnottSmith1994}
M. Knott and C. S. Smith,
``On a Generalization of Cyclic-Monotonicity and Distances among Random
Vectors,''
\emph{Linear Algebra and its Applications} 199 (1994), 363--371.
\url{https://doi.org/10.1016/0024-3795(94)90359-X}

\bibitem{RuschendorfUckelmann2002}
L. Ruschendorf and L. Uckelmann,
``On the $n$-Coupling Problem,''
\emph{Journal of Multivariate Analysis} 81(2) (2002), 242--258.
\url{https://doi.org/10.1006/jmva.2001.2005}

\bibitem{AlvarezEsteban2016}
P. C. Alvarez-Esteban, E. del Barrio, J. A. Cuesta-Albertos, and C. Matran,
``A Fixed-Point Approach to Barycenters in Wasserstein Space,''
\emph{Journal of Mathematical Analysis and Applications} 441(2) (2016),
744--762.
\url{https://doi.org/10.1016/j.jmaa.2016.04.045}

\bibitem{BhatiaJainLim2019}
R. Bhatia, T. Jain, and Y. Lim,
``On the Bures--Wasserstein Distance between Positive Definite Matrices,''
\emph{Expositiones Mathematicae} 37(2) (2019), 165--191.
\url{https://doi.org/10.1016/j.exmath.2018.01.002}

\bibitem{ZemelPanaretos2019}
Y. Zemel and V. M. Panaretos,
``Fr\'echet Means and Procrustes Analysis in Wasserstein Space,''
\emph{Bernoulli} 25(2) (2019), 932--976.
\url{https://doi.org/10.3150/17-BEJ1009}

\bibitem{ChewiEtAl2020}
S. Chewi, T. Maunu, P. Rigollet, and A. J. Stromme,
``Gradient Descent Algorithms for Bures--Wasserstein Barycenters,''
\emph{Proceedings of COLT}, PMLR 125 (2020), 1276--1304.
\url{https://proceedings.mlr.press/v125/chewi20a.html}

\bibitem{AltschulerEtAl2021}
J. M. Altschuler, S. Chewi, P. Gerber, and A. J. Stromme,
``Averaging on the Bures--Wasserstein Manifold: Dimension-Free Convergence of
Gradient Descent,''
\emph{Advances in Neural Information Processing Systems} 34 (2021).
Corrected version: \url{https://arxiv.org/abs/2106.08502}

\bibitem{BrahmachariEtAl2025}
S. Brahmachari, R. Rubboli, and M. Tomamichel,
``A Fixed-Point Algorithm for Matrix Projections with Applications in Quantum
Information,''
\emph{Mathematical Programming} (2025).
\url{https://doi.org/10.1007/s10107-025-02318-w}

\bibitem{MallastoGerolinMinh2022}
A. Mallasto, A. Gerolin, and H. Q. Minh,
``Entropy-Regularized 2-Wasserstein Distance between Gaussian Measures,''
\emph{Information Geometry} 5 (2022), 289--323.
\url{https://doi.org/10.1007/s41884-021-00052-8}

\bibitem{JanatiEtAl2020UnbalancedGaussian}
H. Janati, B. Muzellec, G. Peyre, and M. Cuturi,
``Entropic Optimal Transport between Unbalanced Gaussian Measures Has a
Closed Form,''
\emph{Advances in Neural Information Processing Systems} 33 (2020).
\url{https://proceedings.neurips.cc/paper_files/paper/2020/file/766e428d1e232bbdd58664b41346196c-Paper.pdf}

\bibitem{Chizat2025}
L. Chizat,
``Doubly Regularized Entropic Wasserstein Barycenters,''
\emph{Foundations of Computational Mathematics} (2025).
\url{https://doi.org/10.1007/s10208-025-09724-8}

\bibitem{Minh2023}
H. Q. Minh,
``Entropic Regularization of Wasserstein Distance between Infinite-Dimensional
Gaussian Measures and Gaussian Processes,''
\emph{Journal of Theoretical Probability} 36 (2023), 201--296.
\url{https://arxiv.org/abs/2011.07489}

\bibitem{JanatiCuturiGramfort2020}
H. Janati, M. Cuturi, and A. Gramfort,
``Debiased Sinkhorn Barycenters,''
\emph{Proceedings of ICML}, PMLR 119 (2020).
\url{https://proceedings.mlr.press/v119/janati20a.html}

\bibitem{CumingsMenonShin2020}
R. Cumings-Menon and M. Shin,
``Probability Forecast Combination via Entropy Regularized Wasserstein
Distance,''
\emph{Entropy} 22(9) (2020), 929.
\url{https://doi.org/10.3390/e22090929}

\bibitem{CuturiDoucet2014}
M. Cuturi and A. Doucet,
``Fast Computation of Wasserstein Barycenters,''
\emph{Proceedings of ICML}, PMLR 32 (2014).
\url{https://proceedings.mlr.press/v32/cuturi14.html}

\bibitem{BenamouEtAl2015}
J.-D. Benamou, G. Carlier, M. Cuturi, L. Nenna, and G. Peyre,
``Iterative Bregman Projections for Regularized Transportation Problems,''
\emph{SIAM Journal on Scientific Computing} 37(2) (2015), A1111--A1138.
\url{https://arxiv.org/abs/1412.5154}

\bibitem{LuiseEtAl2019}
G. Luise, S. Salzo, M. Pontil, and C. Ciliberto,
``Sinkhorn Barycenters with Free Support via Frank--Wolfe Algorithm,''
\emph{Advances in Neural Information Processing Systems} 32 (2019).
\url{https://arxiv.org/abs/1905.13194}

\bibitem{ShenEtAl2020}
Z. Shen, Z. Wang, A. Ribeiro, and H. Hassani,
``Sinkhorn Barycenter via Functional Gradient Descent,''
\emph{Advances in Neural Information Processing Systems} 33 (2020).
\url{https://arxiv.org/abs/2007.10449}

\bibitem{LieroMielkeSavare2018}
M. Liero, A. Mielke, and G. Savare,
``Optimal Entropy-Transport Problems and a New Hellinger--Kantorovich Distance
between Positive Measures,''
\emph{Inventiones Mathematicae} 211 (2018), 969--1117.
\url{https://doi.org/10.1007/s00222-017-0759-8}

\bibitem{ChizatEtAl2018}
L. Chizat, G. Peyre, B. Schmitzer, and F.-X. Vialard,
``Scaling Algorithms for Unbalanced Optimal Transport Problems,''
\emph{Mathematics of Computation} 87 (2018), 2563--2609.
\url{https://doi.org/10.1090/mcom/3303}

\bibitem{FrieseckeMatthesSchmitzer2021}
G. Friesecke, D. Matthes, and B. Schmitzer,
``Barycenters for the Hellinger--Kantorovich Distance over $\R^d$,''
\emph{SIAM Journal on Mathematical Analysis} 53 (2021), 62--110.
\url{https://doi.org/10.1137/20M1315555}

\bibitem{ChungPhung2021}
N.-P. Chung and H. M. Phung,
``Barycenters in the Hellinger--Kantorovich Space,''
\emph{Applied Mathematics and Optimization} 84 (2021), 1791--1820.
\url{https://doi.org/10.1007/s00245-020-09695-y}

\bibitem{LieroMielkeTseZhu2025}
M. Liero, A. Mielke, O. Tse, and J.-J. Zhu,
``Evolution of Gaussians in the Hellinger--Kantorovich--Boltzmann Gradient
Flow,'' 2025.
\url{https://doi.org/10.3934/cpaa.2025105}

\bibitem{YangZhang2026}
J. Yang and Y. Zhang,
``Closed Forms for Gaussian KL Unbalanced Optimal Transport without Coupling
Entropy,'' 2026.
\url{https://arxiv.org/abs/2605.02497}

\bibitem{NguyenEtAl2026}
N.-H. Nguyen, L. Q. Dung, H.-P. Nguyen, T. Pham, and N. Ho,
``On Barycenter Computation: Analyzing Semi-Unbalanced Optimal
Transport-Based Method on the Bures--Wasserstein Manifold,''
\emph{AISTATS} (2026).
\url{https://openreview.net/forum?id=EWl46BVj24}

\end{thebibliography}

\end{document}
